% GENERATED FILE. Edit canonical lessons, then run npm run generate:books.
% canonical-source-hash: fnv1a64:f80abead151aff81
% canonical-lessons: TE-C259-raddi, TE-C259-mukhyamaina, TE-C259-sonta, TE-C259-biji, TE-C259-tappanisari

\chapter{Crowded, Important, Own, Busy, Compulsory}
\label{ch:te-crowded-important-own-busy-compulsory}
\hlchaptermodality{259}

\begin{chapteropening}
\textbf{By the end of this chapter:} \emph{I can say crowded, important, own, busy and compulsory.}

Complete the last lesson of chapter 259: I can say crowded, important, own, busy and compulsory.
\end{chapteropening}

\section[\texorpdfstring{\te{రద్దీ} (raddī)}{రద్దీ (raddī)}]{\te{రద్దీ} (raddī) — crowded; a crowd, rush}

\label{lesson:TE-C259-raddi}



\begin{quote}
Before the new one: say the Telugu for development, then the Telugu for a situation.
\end{quote}

\subsection*{You'll want to know: \te{రద్దీ}}
\textbf{\te{రద్దీ}} — \emph{raddī} — \textquotedblleft{}crowded; a crowd, rush\textquotedblright{}.

\begin{grammarlens}[title={What you've built}]
One more word for this chapter.
\end{grammarlens}

\subsection*{Your turn}
\begin{itemize}
\raggedright
  \item \emph{Say it:} \emph{raddī}
  \item \emph{Say it:} \emph{raddī}, once more
  \item \emph{Say it:} say \emph{raddī}
  \item \emph{Recall:} say the Telugu for development, then the Telugu for a situation, then say \emph{raddī} again
\end{itemize}

\begin{tcolorbox}[breakable,colback=teal!4,colframe=teal!35!black,title={Before you move on}]
What does \te{రద్దీ} mean? (\textquotedblleft{}Crowded; a crowd, rush\textquotedblright{}.) Say it once more.
\end{tcolorbox}
\section[\texorpdfstring{\te{ముఖ్యమైన} (mukhyamaina)}{ముఖ్యమైన (mukhyamaina)}]{\te{ముఖ్యమైన} (mukhyamaina) — important}

\label{lesson:TE-C259-mukhyamaina}



\begin{quote}
Before the new one: say the Telugu for a situation, then the Telugu for crowded.
\end{quote}

\subsection*{You'll want to know: \te{ముఖ్యమైన}}
\textbf{\te{ముఖ్యమైన}} — \emph{mukhyamaina} — \textquotedblleft{}important\textquotedblright{}.

\begin{grammarlens}[title={What you've built}]
One more word for this chapter.
\end{grammarlens}

\subsection*{Your turn}
\begin{itemize}
\raggedright
  \item \emph{Say it:} \emph{mukhyamaina}
  \item \emph{Say it:} \emph{mukhyamaina}, once more
  \item \emph{Say it:} say \emph{mukhyamaina}
  \item \emph{Recall:} say the Telugu for a situation, then the Telugu for crowded, then say \emph{mukhyamaina} again
\end{itemize}

\begin{tcolorbox}[breakable,colback=teal!4,colframe=teal!35!black,title={Before you move on}]
What does \te{ముఖ్యమైన} mean? (\textquotedblleft{}Important\textquotedblright{}.) Say it once more.
\end{tcolorbox}
\section[\texorpdfstring{\te{సొంత} (sonta)}{సొంత (sonta)}]{\te{సొంత} (sonta) — own}

\label{lesson:TE-C259-sonta}



\begin{quote}
Before the new one: say the Telugu for crowded, then the Telugu for important.
\end{quote}

\subsection*{You'll want to know: \te{సొంత}}
\textbf{\te{సొంత}} — \emph{sonta} — \textquotedblleft{}own\textquotedblright{}.

\begin{grammarlens}[title={What you've built}]
One more word for this chapter.
\end{grammarlens}

\subsection*{Your turn}
\begin{itemize}
\raggedright
  \item \emph{Say it:} \emph{sonta}
  \item \emph{Say it:} \emph{sonta}, once more
  \item \emph{Say it:} say \emph{sonta}
  \item \emph{Recall:} say the Telugu for crowded, then the Telugu for important, then say \emph{sonta} again
\end{itemize}

\begin{tcolorbox}[breakable,colback=teal!4,colframe=teal!35!black,title={Before you move on}]
What does \te{సొంత} mean? (\textquotedblleft{}Own\textquotedblright{}.) Say it once more.
\end{tcolorbox}
\section[\texorpdfstring{\te{బిజీ} (bijī)}{బిజీ (bijī)}]{\te{బిజీ} (bijī) — busy}

\label{lesson:TE-C259-biji}



\begin{quote}
Before the new one: say the Telugu for important, then the Telugu for own.
\end{quote}

\subsection*{You'll want to know: \te{బిజీ}}
\textbf{\te{బిజీ}} — \emph{bijī} — \textquotedblleft{}busy\textquotedblright{}.

\begin{grammarlens}[title={What you've built}]
One more word for this chapter.
\end{grammarlens}

\subsection*{Your turn}
\begin{itemize}
\raggedright
  \item \emph{Say it:} \emph{bijī}
  \item \emph{Say it:} \emph{bijī}, once more
  \item \emph{Say it:} say \emph{bijī}
  \item \emph{Recall:} say the Telugu for important, then the Telugu for own, then say \emph{bijī} again
\end{itemize}

\begin{tcolorbox}[breakable,colback=teal!4,colframe=teal!35!black,title={Before you move on}]
What does \te{బిజీ} mean? (\textquotedblleft{}Busy\textquotedblright{}.) Say it once more.
\end{tcolorbox}
\section[\texorpdfstring{\te{తప్పనిసరి} (tappanisari)}{తప్పనిసరి (tappanisari)}]{\te{తప్పనిసరి} (tappanisari) — compulsory}

\label{lesson:TE-C259-tappanisari}



\begin{quote}
Before the new one: say the Telugu for own, then the Telugu for busy.
\end{quote}

\subsection*{You'll want to know: \te{తప్పనిసరి}}
\textbf{\te{తప్పనిసరి}} — \emph{tappanisari} — \textquotedblleft{}compulsory\textquotedblright{}.

\begin{grammarlens}[title={What you've built}]
One more word for this chapter.
\end{grammarlens}

\subsection*{Your turn}
\begin{itemize}
\raggedright
  \item \emph{Say it:} \emph{tappanisari}
  \item \emph{Say it:} \emph{tappanisari}, once more
  \item \emph{Say it:} say \emph{tappanisari}
  \item \emph{Recall:} say the Telugu for own, then the Telugu for busy, then say \emph{tappanisari} again
\end{itemize}

\begin{tcolorbox}[breakable,colback=teal!4,colframe=teal!35!black,title={Before you move on}]
What does \te{తప్పనిసరి} mean? (\textquotedblleft{}Compulsory\textquotedblright{}.) Say it once more.
\end{tcolorbox}
